\documentclass[11pt]{article}
\usepackage{amsmath,amssymb,amsthm,hyperref}
\newtheorem{theorem}{Theorem}\newtheorem{proposition}[theorem]{Proposition}
\title{Controlled native joint crossing and the exact temporal factorization boundary}
\date{16 November 2026}
\begin{document}\maketitle
\paragraph{Status.} The finite constant-polynomial construction, native comparisons and causal factorization criterion have passed independent written review. No proof-assistant certification or TeX compilation is claimed. This isolates the original polynomial crossing interface where arbitrary variable arities are harmless and the context-multiplication interface where an unrestricted all-coupling law fails. It supplies a constructive controlled crossing, full native grouping maps, and an exact nonanticipation criterion. A complete recursive strategy/compiler theorem is a separate integration step.

\section{Original inputs and explicit controls}
Use the original O1 context rig $\mathbb R_{\ge0}^{\mathrm{disc}}\times\mathcal A$, where $\mathcal A$ has cocartesian addition and canonical distributors, and $\mathcal E=\mathcal A/I$ has finite products. A row consists of finite raw witnesses $W_b$, normalized law $\mu_b$, and actual augmented coefficient objects $C_{b,w}\in\mathcal E$. Native maps retain all actual witness and coefficient maps. This is the full native original augmented-family category for this source, by the coproduct universal property.

For clarity the first construction uses original constant polynomial content with arbitrary finite arities: a global collection $Q=(Q_n)_{n\ge0}$, actual arity-preserving shape maps, and extension $F_QB=\coprod_n Q_n\times B^n$. All unused shapes and their maps remain in $Q$. Coefficient-decorated content can be included with its actual original maps; the context-side theorem does not infer those maps from their carriers.

At a current content point $y=(o,b_1,\ldots,b_n)$, the occurrence positions $[n]$ and their input labels are known before the continuation context is sampled. Define
\[
 W_y=\prod_{i=1}^n W_{b_i},\qquad
 \mathsf{Coup}(y)=\{q\in\mathsf D(W_y): (\pi_i)_*q=\mu_{b_i}\text{ for every }i\}.
\]
This is nonempty, since the independent product law is one member. Empty occurrence families have singleton witness and its unique law. The control interface is the actual set
\[
 \widetilde Y=\{(y,q):q\in\mathsf{Coup}(y)\}\longrightarrow F_QB.
\]
At $(y,q)$ use the native context row with raw set $W_y$, law $q$, and coefficient $\prod_i C_{b_i,w_i}$. Its body fills the original ordered polynomial positions with these witnesses, exactly as in the old display crossing. The new control is not a probability mixture, not a selected law hidden in the grade, and not an assumed descent of a policy through an erased input. It is a supplied pure input. The contextual extension fibre remains finite even if the control set is infinite.

\begin{theorem}[Controlled joint crossing and native transport]
The displayed construction is functorial on original context maps, original arity-preserving content maps, and compatible pure endpoint squares. A context comparison with row witness maps $r_b$ sends $(y,q)$ to $(y,(\prod_i r_{b_i})_*q)$. A content map sends $y$ to its target shape with the same arity and ordered labels, leaving the joint control on those witness coordinates unchanged. The actual context coefficient map is the product of the primitive coefficient maps; original content transport/display maps supply the independent grade component. No whole-system joint square is assumed.
\end{theorem}
\begin{proof}
Projection commutation gives the required target marginal at every occurrence. Product universality constructs the coefficient map and proves its identity/composition laws on actual arrows. The original polynomial body equation follows on its shape and each ordered output coordinate; its grade map remains the supplied original content arrow, not a carrier quotient. Pure reindexing merely substitutes the input point, and output maps act at each coordinate. For a general original content input satisfying O1's display/cartesian hypotheses, the same body is obtained from its display bijection and its grade arrow by the original cartesian/opcartesian factorization. Its existence is conditional on those original hypotheses, not on any new global feedback law. This observation gives a conditional display-body construction only; no full controlled contentad or multiplication theorem for arbitrary L1--L6 inputs is inferred from it. The proved grouping and temporal statements below use the explicitly specified constant polynomial source. A fixed policy $y\mapsto q_y$ induces an ordinary process only after it is supplied; transporting such a policy to a smaller visible state still needs descent. Keeping $(y,q)$ as the interface makes the open construction functorial without asserting that descent.
\end{proof}

\section{Variable-shape grouping is exact at a fixed input}
Suppose a current content point supplies a finite addressed two-level tree, with outer positions $i$ and inner positions $j\in J_i$. All these positions are part of that input, before the context draw. A flat law $q$ on $\prod_i\prod_jW_{ij}$ determines its group marginals $q_i$ and the same law on the grouped tuple. Conversely a coupling of group tuples with specified marginals $q_i$, each itself coupling the required primitive marginals, determines one flat law. These assignments are inverse, including nullary groups. Actual coefficient products have the canonical isomorphism
\[
 \prod_i\prod_j C_{ij}\cong\prod_{(i,j)}C_{ij}.
\]
The original content substitution retains its full outer/inner shape tuple and its original native grade comparisons. Triple regrouping agrees on the complete joint law and every actual coefficient projection, proving native coherence. Thus the content-multiplication grouping itself does not require fixed global arity: what matters here is that the current addressed tree is already an input. In the constant ordered-polynomial source, native shape maps preserve the ordered positions identically; the displayed regrouping uses the canonical addressed tuple bijection. No additional arbitrary position permutation is attributed to those native grade maps. Unused global shapes are not discarded.

\section{The context-multiplication obstruction}
At one fixed current input, let $i\in I$ index simultaneous occurrences. Their first-stage witness laws are $\mu_i$ on $W_i$; their second-stage laws are $\nu_{i,w_i}$ on $Z_{i,w_i}$. A locally bound row has law
\[
 \beta_i(w_i,z_i)=\mu_i(w_i)\nu_{i,w_i}(z_i).
\]
An arbitrary joint law $q$ coupling the $\beta_i$ need not admit a sequential interpretation in which the second-stage primitive marginals remain correct after the complete first-stage witness vector is known.

Take two occurrences and four bits $(X_1,Y_1,X_2,Y_2)$. Choose $X_1,X_2$ independently fair and set $Y_1=X_2$, $Y_2=X_1$. Each local pair $(X_i,Y_i)$ has the independent fair law, so $q$ is an admitted coupling of the two locally bound rows. But after observing the complete first-stage vector $(X_1,X_2)$, both second-stage outputs are deterministic. For example $\Pr(Y_1=1\mid X_1=0,X_2=1)=1$, not $1/2$. Thus an unrestricted inversion from whole-subtrajectory coupling to fresh second-stage joint sampling is impossible. This is not a counterexample to the original independent law or to a directed sequential-to-joint comparison.

\begin{theorem}[Exact causal factorization and its retained zero-row fibre]
Let $q$ be a normalized joint law on the complete pairs whose $i$th pair marginal is $\beta_i(w_i,z_i)=\mu_i(w_i)\nu_{i,w_i}(z_i)$ for every $i$, and let $p(\mathbf w)$ be its first-stage marginal. It is generated by a first-stage coupling $p$ and a family of second-stage couplings $r_{\mathbf w}$ with marginals $\nu_{i,w_i}$ through
\[
 q(\mathbf w,\mathbf z)=p(\mathbf w)r_{\mathbf w}(\mathbf z) \tag{F}
\]
if and only if, for every $\mathbf w,i,z_i$,
\[
 \sum_{\mathbf z:\,\mathbf z_i=z_i}q(\mathbf w,\mathbf z)
 =p(\mathbf w)\nu_{i,w_i}(z_i).\tag{N}
\]
At positive-mass $\mathbf w$ the factor $r_{\mathbf w}$ is uniquely $q(\mathbf w,-)/p(\mathbf w)$. At zero-mass rows it is unconstrained except for its specified primitive marginals. Hence flattening a strategy that retains all zero-row conditional choices to its whole law is generally noninjective; no such choices are silently identified.
\end{theorem}
\begin{proof}
Equation (F) and the marginal conditions give (N). Conversely divide at positive rows. At zero rows every $q(\mathbf w,\mathbf z)$ is zero by nonnegativity, so choose any coupling of the normalized second-stage marginals, for example their independent product. The resulting $r_{\mathbf w}$ gives (F). The first-stage marginal $p$ already has marginals $\mu_i$ because $q$ couples the local bound rows. The stated fibre description follows directly. It is not a claim of a natural canonical choice at zero rows.
\end{proof}

\section{Actual native directed multiplication}
A sequential choice consists of the first coupling $p$ and every conditional coupling $r_{\mathbf w}$, including zero-weight rows. Flattening sends it to (F), an admitted coupling of the locally bound rows. If first and second context coefficients are $C_{i,w_i}$ and $D_{i,w_i,z_i}$, sequential joint assembly has coefficient
\[
 (\prod_i C_{i,w_i})\otimes(\prod_i D_{i,w_i,z_i}),
\]
while one-shot joint assembly of the bound rows has $\prod_i(C_{i,w_i}\otimes D_{i,w_i,z_i})$. The canonical arrow from the former to the latter is obtained by tensoring the two projections to each $i$ and invoking target product universality. Its naturality is a native coefficient equation; it need not be invertible.

For three sequential stages with full retained prefix controls, both temporal bracketings evaluate to the same probability product
\[
 p(\mathbf w)r_{\mathbf w}(\mathbf z)t_{\mathbf w,\mathbf z}(\mathbf v)
\]
on the same complete addressed history, and both coefficient arrows have the same projected primitive tensors. Thus the directed evaluation comparison is coherent when the prefix controls remain part of the strategy. This does not turn arbitrary coupling of completed subtrajectory laws into a causal strategy: the preceding counterexample fails (N).

The staged coefficient convention in this section is a tensor across time blocks of products within each simultaneous frontier. It is not identical to a strategy evaluator using categorical product across all executed nodes. There is a canonical directed map to that all-node product: augment the other temporal blocks and then project to the selected node, and assemble these arrows by product universality. For any finite sequence of stages, all cut/reassociation paths have the same node projections, so the comparison is coherent. No tensor/product isomorphism is claimed; this is an explicit bridge between two coefficient doctrines.

A new conditional choice may depend on the full already revealed routing state. Its primitive marginal remains the specified row law at that state; correlations among the simultaneous new outputs are unrestricted. For example both second-stage fair outputs may use one shared fair bit, and the selected diagonal versus antidiagonal coupling may depend on the entire first-stage history. This is genuine shared interference, not a return to independent tensor. A common random scope reused after partial observation must be included in state with its actual conditional law; it cannot be treated as a new unconditioned fair draw.

\paragraph{Remaining integration obligation.} A complete strategy theorem must define finite-round/open-history strategy objects, their binary substitution retaining the routing state and all conditional controls, evaluation into actual native grades, and a cut theorem preserving the schedule. It must relate finite terminating histories to the variable-arity recursive grades without treating those grades as homogeneous. This paper supplies the source-native crossing, the exact temporal obstruction, and the conditional multiplication comparison needed for that theorem; it does not claim the whole compiler or feedback integration by assertion.
\end{document}
